\begin{abstract}
围绕“过期意大利面-42号混凝土”这一趣味题目，本文构造低碳配合比优化的学术演示。
“42号”仅为案例编号，所有材料参数、性能响应及计算结果均属于合成情景，没有采用实际材料试验数据。
研究旨在展示废弃物添加方案的建模、求解与解释过程，不将废弃物利用预设为减碳收益。

以每立方米模型拌合物为功能单位，选取胶凝材料总量、矿渣比例、有效水胶比和干意面用量为决策变量，
建立含吸水补偿的绝对体积配料模型。在假设强度、坍落度、总水量与砂量约束下，
同时最小化材料碳排放和成本、最大化合成强度响应。
采用约束支配的 NSGA-II 求取近似 Pareto 集，并以固定归一化下的等权理想点距离选取折中方案。
设置五个随机种子和同评价预算的均匀随机搜索，检查计算重复性及候选集质量。

主运行得到 \FrontCount{} 个可行非支配方案。折中方案的碳排指标为
\CompCarbon{} kg CO$_2$e/m$^3$，材料成本为 \CompCost{} 元/m$^3$，
合成强度为 \CompStrength{} MPa；相较预设基准，碳排与成本分别降低
\CarbonReduction{}\% 和 \CostReduction{}\%，强度减少 \StrengthLoss{} MPa。
折中方案干意面用量仅为 \CompPasta{} kg/m$^3$，表明当前响应与核算假设倾向低掺量，
结果不支持“加入意面必然提高性能或降低碳排”的结论。
固定配合比的扰动检验进一步说明：强度余量与坍落度余量必须分别检查，
有限搜索和单因素分析均不能替代材料试验与重新标定。
本文给出完整参数、可运行程序和自动生成的图表，形成可追溯的多目标优化演示。

\keywords{NSGA-II；配合比优化；合成数据；Pareto 权衡；碳排核算}
\end{abstract}
